%-------------------------
% Resume in Latex
% Author : Shashank Sharma
%-------------------------

\documentclass[letterpaper,10pt]{article}

% Packages
\usepackage[top=0.4in, bottom=0.4in, left=0.5in, right=0.5in]{geometry}
\usepackage{titlesec}
\usepackage{enumitem}
\usepackage{hyperref}
\usepackage{microtype}
\usepackage{needspace}

% Links
\hypersetup{
  colorlinks,
  urlcolor=blue,
  linkcolor=blue
}

% Page layout
\pagestyle{empty}
\urlstyle{same}
\raggedbottom
\raggedright
\setlength{\tabcolsep}{0pt}

% Section formatting
\titleformat{\section}
  {\scshape\raggedright\large}
  {}
  {0em}
  {}
  [\titlerule]

\titlespacing*{\section}
  {0pt}
  {3pt}
  {2pt}

% Bullet lists
\setlist[itemize]{
  leftmargin=14pt,
  itemsep=0pt,
  topsep=0pt,
  parsep=0pt,
  partopsep=0pt
}

% Resume headings
\newcommand{\resumeHeadingTwoLine}[4]{%
  \needspace{3\baselineskip}%
  \begin{tabular*}{\textwidth}{l@{\extracolsep{\fill}}r}
    \textbf{#1} & #2 \\
    \textit{\small #3} & \textit{\small #4}
  \end{tabular*}%
}

\newcommand{\resumeHeadingOneLine}[2]{%
  \begin{tabular*}{\textwidth}{l@{\extracolsep{\fill}}r}
    \textbf{#1} & \textit{\small #2}
  \end{tabular*}%
}

\newcommand{\resumeSubheading}[2]{%
  \begin{tabular*}{\textwidth}{l@{\extracolsep{\fill}}r}
    \textit{\small #1} & \textit{\small #2}
  \end{tabular*}%
}

\newcommand{\resumeItem}[1]{\item\small #1}

%-------------------------------------------
%%%%%%  CV STARTS HERE  %%%%%%%%%%%%%%%%%%%%%%%%%%%%
\begin{document}



%----------HEADING-----------------
\begin{tabular*}{\textwidth}{l@{\extracolsep{\fill}}c@{\extracolsep{\fill}}r}
    \href{https://www.linkedin.com/in/ssharma1991/}{linkedin.com/in/ssharma1991} &
    \textbf{\LARGE Shashank Sharma} &
    \href{mailto:shashank.sharma.1991@outlook.com}{shashank.sharma.1991@outlook.com}\\
    \href{http://sharmashashank.com/}{sharmashashank.com} &
    Fremont, California \quad 631-512-0029 &
    \href{https://github.com/ssharma1991}{github.com/ssharma1991}\\
\end{tabular*}



%-----------EXPERIENCE-----------------
\section{Experience}

\resumeHeadingTwoLine
{Lucid Motors}{Newark, CA}
{Senior Software Engineer (Maps and Navigation)}{Nov 2023 -- Present}
\begin{itemize}
\resumeItem{Led development of HD map software on Nvidia Xavier/Orin with QNX, integrating GNSS sensor data with TomTom APIs.}
\resumeItem{Delivered ADAS features (Lane Change Assist, Adaptive Curve Speed Control, Hands-Free Highway Assist) through cross-functional collaboration with platform, perception, motion planning, cloud, and field testing teams.}
\resumeItem{Enhanced system stability and runtime reliability through memory, CPU, and bandwidth profiling, improving offline performance and reducing error rates.}
\resumeItem{Created interactive visualization and debugging tools with Folium, Plotly, and QGIS, enhancing map data analysis and accelerating issue resolution for ADAS development.}
\end{itemize}

\resumeHeadingTwoLine
{Dematic, Kion Mobile Automation}{Holland, MI}
{Machine Learning Engineer (Perception, Localization, and Mapping)}{Sept 2020 -- Oct 2023 }
\begin{itemize}
\resumeItem{Developed onboard mapping and localization software for autonomous guided vehicles (AGVs) in warehouse environments.}
\resumeItem{Achieved sub‑centimeter localization accuracy by refining reflector and line‑extraction algorithms.}
\resumeItem{Accelerated localization pipeline by 30\% through optimizing feature-based association algorithms.}
\resumeItem{Reduced technical debt in the mapping codebase by simplifying APIs, decoupling modules, and introducing caching threads; enhanced FastSLAM implementation for robust real‑time performance.}
\resumeItem{Developed YOLOv3-based perception model for real-time detection of persons, forklifts, and pallets using 3D camera data.}
\end{itemize}

\resumeHeadingTwoLine
{Stony Brook University}{Stony Brook, NY}
{Research Assistant}{May 2017 -- Aug 2020}
\begin{itemize}
\resumeItem{Developed machine‑learning and algebraic algorithms for simulating and synthesizing robotic mechanisms, leading to multiple peer‑reviewed ASME publications.}
\resumeItem{Built  \href{http://cadcam.eng.sunysb.edu/}{MotionGen}, a web-based mechanism design framework using the MEAN stack, delivering a RESTful service with cross‑platform iOS/Android apps via Apache Cordova.}
\end{itemize}
	


%-----------EDUCATION-----------------
\section{Education}
\resumeHeadingTwoLine
{Stony Brook University}{Stony Brook, NY}
{Ph.D. in Mechanical Engineering (Design and Robotics; Minor: Applied Mathematics), GPA 3.95}{Aug 2015 -- Aug 2020}
\begin{itemize}
\resumeItem{\textbf{Relevant Courses:} Robotics, Dynamics, Control, Kinematics, Stress Analysis, Optimization, Algorithms}
\end{itemize}



%-----------PROJECTS-----------------
\section{Selected Projects}

% Mapping & Localization
\resumeSubheading{Mapping \& Localization}{}
\begin{itemize}
\resumeItem{Developed SLAM pipelines (gmapping, RTAB‑Map) with LiDAR and 3D cameras, generating occupancy grids and loop‑closed maps for robust AGV/AMR localization.}
\resumeItem{Implemented particle filter and extended Kalman filter for multi-sensor fusion (GPS, LiDAR, RADAR), reducing localization drift under noisy sensor conditions and improving reliability of autonomous navigation pipelines.}
\end{itemize}

% ML & Perception
\resumeSubheading{Perception \& Machine Learning}{}
\begin{itemize}
\resumeItem{Built a stereo‑camera distance estimator for front‑view vehicles. Used SGBM disparity calculation for depth and YOLO for detection, enabling real‑time ADAS perception.}
\resumeItem{Implemented 3D reconstruction with structure‑from‑motion (SfM). Applied SIFT feature matching, computed fundamental and essential matrices, aligned image sequences, and generated 3D point‑cloud models from monocular data.}
\resumeItem{Developed a 3D object‑detection framework combining camera and LiDAR. Applied sensor fusion to project point clouds into images, detect objects, and estimated 3D bounding boxes, improving perception accuracy in autonomous driving.}
\resumeItem{Built a traffic‑sign classification pipeline using CNNs trained on 40+ classes in Keras, achieving over 95\% accuracy.}
\resumeItem{Created real-time lane detection pipeline using BEV transforms and polynomial fitting.}
\end{itemize}

% Planning & Control
\resumeSubheading{Planning \& Control}{}
\begin{itemize}
\resumeItem{Designed an FSM trajectory planner and PID controller for autonomous highway driving; integrated perception, localization, planning, and control modules within ROS to achieve end‑to‑end autonomy.}
\resumeItem{Implemented autonomous pick-and-drop operations in Gazebo simulations by combining AMCL-based localization with path planning algorithms, demonstrating full autonomy in simulated warehouse workflows.}
\end{itemize}



%-------- SKILLS------------
\section{Technical Proficiency}
	\begin{itemize}
		\resumeItem{\textbf{Programming:} C\texttt{++}, Python, JavaScript, Matlab, Mathematica, Delphi}
		 \resumeItem{\textbf{Machine Learning \& Robotics:} TensorFlow, PyTorch, Keras, OpenCV, Scikit-learn, Pandas, Open3D, CNNs}
		 \resumeItem{\textbf{Mapping \& Localization:} ROS, Gazebo, Rviz, gmapping, RTAB-Map, AMCL, EKF, particle filter, QGIS}
		\resumeItem{\textbf{Tools \& Infrastructure:} Git, Anaconda, Jupyter, VirtualBox, Optimization(LTTng, Perf, Valgrind)}
	\end{itemize}


%-------- RESEARCH PAPERS------------
\section{Selected Publications}
\begin{itemize}
\resumeItem{Sharma S., Purwar A.; \textbf{A Machine Learning Approach to Solve the Alt–Burmester Problem for Synthesis of Defect-Free Spatial Mechanisms.} ASME J. Computing and Information Science in Engineering; doi:10.1115/1.4051913}
\resumeItem{Sharma S., Purwar A.; \textbf{Path Synthesis of Defect-Free Spatial 5-SS Mechanisms Using Machine Learning.}, ASME IDETC-CIE2020; doi:10.1115/DETC2020-22731}
\resumeItem{Sharma S., Purwar A., Ge Q.J.; \textbf{Motion Synthesis Approach to Solving Alt-Burmester Problem by Exploiting Fourier Descriptor Relationship Between Path and Orientation.}, ASME J. Mechanisms Robotics; doi:10.1115/1.4042054}
%\resumeItem{Sharma S., Purwar A., Ge Q.J.; \textbf{An Optimal Parametrization Scheme for Path Generation Using Fourier Descriptors for Four-Bar Mechanism Synthesis.}, ASME J. Computing and Information Science in Engineering; doi:10.1115/1.4041566}
\end{itemize}

\end{document}
